% !TEX root = ../tfg_etsiinf_NombreApellidos.tex
\chapter*{Resumen} \label{chp:abstract}

Este trabajo desarrolla un entorno de aprendizaje por refuerzo para drones
sobre el simulador dinámico de multirrotores de Aerostack2, un \emph{framework}
de robótica aérea basado en ROS 2. Sobre este entorno se entrena y evalúa un
controlador de velocidad de alto nivel. La tarea consiste en alcanzar objetivos
con posición y orientación aleatorias desde estados iniciales también
aleatorios, mediante observaciones relativas normalizadas y comandos continuos
de velocidad lineal y de guiñada.

La primera contribución es la infraestructura experimental. Se implementa un
entorno Gymnasium integrado con Aerostack2, un reinicio determinista del
simulador con certificación de accionabilidad, guards predictivos de acción, un
supervisor para campañas largas con reanudación desde \emph{checkpoint} y
reciclado preventivo, y ejecución vectorizada de dos drones en un mismo
simulador. En conjunto, estos componentes permiten iniciar episodios solo cuando
el vehículo responde físicamente a los comandos, corregir aproximaciones a
fronteras no observables y duplicar el muestreo efectivo.

La segunda contribución es la resolución de la tarea de navegación. Con una
recompensa mínima basada en distancia y un curriculum de tres etapas que amplía
progresivamente el rango de distancias, la tasa de éxito pasa de valores
inferiores al 1\,\% al 100\,\% en una evaluación determinista de 100 episodios
reproducibles por semillas.

La evaluación final compara la política aprendida con un controlador PID de
referencia sobre los mismos episodios. Ambos alcanzan el 100\,\% de éxito, pero
el PID resuelve la tarea nominal con mayor rapidez y mejor precisión de
orientación. Este resultado es coherente con la naturaleza del problema: el
control de velocidad hacia un objetivo estático favorece a una ley proporcional.
Por ello, el valor del enfoque aprendido no reside en superar al PID en este
caso concreto, sino en alcanzar una paridad funcional y en proporcionar una
infraestructura aplicable a tareas sin una ley de control evidente.

El trabajo también identifica dos limitaciones de ingeniería del simulador: la
degradación acumulativa asociada a los reinicios de servicio y la persistencia
indefinida de las referencias de velocidad. Estos hallazgos motivan, como línea
futura principal, el entrenamiento sobre \emph{bindings} directos de la dinámica
sin ROS 2 en el bucle.

\textbf{Palabras Clave}: aprendizaje por refuerzo, vehículos aéreos no
tripulados, Aerostack2, ROS 2, PPO, curriculum, simulación, control de
velocidad.

%%%%%%%%%%%%%%%%%%%%%%%%%%%%%%%%%%%%%%%%%%%%%%%%%%%%%%%%%%%%%%%%%%%%%%%%%%%%%%%%

\newpage

%%%%%%%%%%%%%%%%%%%%%%%%%%%%%%%%%%%%%%%%%%%%%%%%%%%%%%%%%%%%%%%%%%%%%%%%%%%%%%%%

\chapter*{Abstract}

This work develops a reinforcement learning environment for drones on top of
the Aerostack2 multirotor dynamics simulator, a ROS 2 based aerial robotics
framework. The environment is used to train and evaluate a high-level velocity
controller. The task consists of reaching targets with random positions and
headings from random initial states, using normalized relative observations and
continuous linear and yaw velocity commands.

The first contribution is the experimental infrastructure. The work implements
a Gymnasium environment integrated with Aerostack2, a deterministic simulator
reset with actionability certification, predictive action guards, a supervisor
for long training campaigns with checkpoint resumption and preventive
recycling, and vectorized execution of two drones in a single simulator.
Together, these components make it possible to start episodes only when the
vehicle physically responds to commands, correct approaches to unobservable
boundaries, and double the effective sampling rate.

The second contribution is solving the navigation task. With a minimal
distance-based reward and a three-stage curriculum that progressively expands
the distance range, the success rate increases from below 1\,\% to 100\,\% in
a deterministic evaluation of 100 seed-reproducible episodes.

The final evaluation compares the learned policy with a reference PID
controller on the same episodes. Both reach 100\,\% success, but the PID solves
the nominal task faster and with better final heading accuracy. This result is
consistent with the nature of the problem: velocity control towards a static
target favors a proportional law. Therefore, the value of the learned approach
does not lie in outperforming the PID in this specific case, but in achieving
functional parity and providing an infrastructure that can be applied to tasks
without an obvious control law.

The work also identifies two engineering limitations of the simulator:
cumulative degradation associated with reset service calls and the indefinite
persistence of velocity references. These findings motivate, as the main future
line, training over direct dynamics bindings without ROS 2 in the loop.

\textbf{Keywords}: reinforcement learning, unmanned aerial vehicles,
Aerostack2, ROS 2, PPO, curriculum learning, simulation, velocity control.
